\section{Comparison and transport}\label{sec:foundations}

The diagonalization argument repeatedly places several copies of a
positive element inside one hereditary subalgebra.  Comparison of
Cuntz classes alone does not record enough information for this use:
the copies must have a common implementation, their products must
remain in the algebra, and their norms must agree in every
representation.  We establish these facts here, together with the
ideal and cutoff rules used to retain norm signals.

Cuntz introduced the comparison relation in \cite{Cuntz1978}.
The proper-target ideal comparison in
Lemma~\ref{lem:comparison-rules}(iii) is the standard consequence
recorded in \cite[Proposition~3.5(ii)]{KR2000}; it is also used with
that reference in
\cite[proof of Proposition~2.38]{KwasniewskiMeyer2023}.
The cutoff implementation below is related to R{\o}rdam's
comparison lemma \cite[Proposition~2.4]{Rordam1992}, in the form
recalled in \cite[Section~2.1]{AntoinePereraThiel2020}.
We give the rectangular and representation statements needed here
with their proofs.

All comparisons below are stable: finite positive matrices are
placed in a common stabilization by padding with zeros.  For a
\(C^*\)-algebra \(D\) and \(h\in D_+\), let \(\Id_D(h)\)
denote the closed two-sided ideal generated by \(h\), and put
\[
 \Her_D(h)=\overline{hDh},\qquad
 e_h=\mathbf 1_{(0,\infty)}(h)\in D^{**}.
\]
We omit the subscript when the ambient algebra is clear.  Every
support projection is taken in the universal bidual.  Rectangular
matrices over \(C^{**}\) are viewed in the corresponding corners of
a matrix algebra over \(C^{**}\).  The normal extension of a
representation sends \(e_h\) to \(e_{\pi(h)}\): apply normality to
\(h(h+1/j)^{-1}\uparrow e_h\).  When using a unit or a support
projection in a degenerate representation, we restrict to its
essential subspace.  This leaves the norm of every represented
algebra element unchanged.

\begin{lemma}[Elementary comparison rules]\label{lem:comparison-rules}
Let \(D\) be a \(C^*\)-algebra.
\begin{enumerate}[label=\textup{(\roman*)}]
\item For \(h\in D_+\),
\[
 \Her_D(h)=D\cap e_hD^{**}e_h.
\]
If \(0\le u\le K v\) for \(u,v\in D_+\) and \(K\ge0\), then
\(u\in\Her(v)\) and \(u\precsim v\).  More generally, every
positive element of \(\Her(v)\) is Cuntz subequivalent to \(v\).
\item For a rectangular algebra matrix \(z\), one has
\(z^*z\CuEq zz^*\).  Cuntz subequivalence is transitive, is
preserved by block sums, and is closed under norm limits in its
first variable when the second variable is fixed.  Every
\(*\)-homomorphism preserves Cuntz subequivalence and equivalence,
after amplification.  For integers \(m,n\ge1\), \(d\in M_n(D)_+\),
and \(v\in M_m(D)_+\), stable comparison \(d\precsim v\) is
equivalent to the existence of \(r_j\in M_{m,n}(D)\) with
\(r_j^*vr_j\to d\) in norm.
\item If \(d,v\) are positive elements in a common stabilization,
\(d\in\Id(v)\), and \(v\) is properly infinite, then
\(d\precsim v\).
\end{enumerate}
\end{lemma}

\begin{proof}
We first prove the bidual description in (i), including its norm
closure assertion.  One inclusion follows on \(hDh\) and then by
closure.  Conversely, suppose \(u\in D_+\) and \(e_hu=u\).
Take increasing continuous positive contractions \(f_j\) with
\(f_j(0)=0\) and \(f_j(t)\uparrow1\) for \(t>0\).  On the compact
state space of the unitization, the continuous functions
\[
 \varphi\longmapsto
 \varphi\bigl(u(1-f_j(h))u\bigr)
\]
decrease pointwise to zero.  Indeed, the normal extension of each
state evaluates the strong limit \(f_j(h)\uparrow e_h\).
Dini's theorem makes this convergence uniform.  The positive norm
formula and \((1-f_j(h))^2\le1-f_j(h)\) then give
\(\|(1-f_j(h))u\|\to0\).  Thus
\(u=\lim_j f_j(h)u f_j(h)\in\Her(h)\).
Here \(f_j(h)\in\Her(h)\) follows from continuous functional
calculus, by first cutting \(f_j\) off near zero.  For a general
element supported on both sides by \(e_h\), apply the positive case
to the positive and negative parts of its real and imaginary parts.

If \(0\le u\le Kv\), compression by \(1-e_v\) gives
\(u^{1/2}(1-e_v)=0\).  Hence \(u\) is supported by \(e_v\) and
belongs to \(\Her(v)\).  If \(u\in\Her(v)_+\), take a cutoff
approximate identity \(g_j(v)\) for \(\Her(v)\), with every \(g_j\)
zero near zero.  The quotient function below is continuous when
extended by zero at zero, and gives elements
\[
 x_j=\frac{g_j(v)}{\sqrt v}\,u^{1/2}\in D,\qquad
 x_j^*v x_j=u^{1/2}g_j(v)^2u^{1/2}\longrightarrow u.
\]
This proves the comparison assertion.

For (ii), the rectangular coefficient
\(z(z^*z+1/j)^{-1/2}\) compares \(z^*z\) to \(zz^*\), since its
conjugation of \(zz^*\) equals
\[
 (z^*z)^2(z^*z+1/j)^{-1}\longrightarrow z^*z.
\]
Apply the same calculation to \(z^*\) for the reverse comparison.
The inverses lie in unitizations and are multiplied by algebra
matrices.  Block-diagonal coefficients prove preservation under
block sums.  For transitivity, first approximate \(u\) by \(x^*vx\),
then approximate \(v\) by \(y^*wy\) within the requested error
divided by \(\max\{1,\|x\|^2\}\); the coefficient \(yx\) compares
\(u\) to \(w\).  A diagonal choice of coefficients proves closure
under norm limits in the first variable.

If \(\phi:D\to E\) is a \(*\)-homomorphism, apply
\(\phi\otimes\id_{\Kc}\) to a comparison sequence; continuity
preserves its norm limit.  Applying this in both directions preserves
equivalence.  For finite \(d,v\), let \(P_n,P_m\) be the scalar
coordinate projections in the multiplier algebra of \(D\otimes\Kc\).
If \(x_j^*vx_j\to d\) in the stabilization, then
\[
 r_j=P_mx_jP_n\in M_{m,n}(D),\qquad
 r_j^*vr_j=P_nx_j^*vx_jP_n\longrightarrow d.
\]
Conversely, rectangular coefficients embed in the stabilization and
give the same convergence there.  The multiplier compressions retain
algebra membership without a unit, and the argument includes zero
matrices.

For (iii), work in the common stabilization \(D\).  The closed span
of \(Dv^{1/2}D\) is \(\Id_D(v)\), so finite sums
\(\sum_{i=1}^r x_i v^{1/2}y_i\) approximate \(d^{1/2}\).
For such a sum, define
\[
 X=[x_1\ \cdots\ x_r]\in M_{1,r}(D),\qquad
 Y=[y_1\ \cdots\ y_r]^T\in M_{r,1}(D),
\]
and put \(V=\operatorname{diag}(v,\ldots,v)\in M_r(D)\).  Then
\[
 z=\sum_{i=1}^r x_i v^{1/2}y_i=X V^{1/2}Y.
\]
The order and rectangular rules already proved give
\[
 z^*z\le\|X\|^2Y^*VY,\qquad
 Y^*VY\CuEq V^{1/2}YY^*V^{1/2}\le\|Y\|^2V.
\]
It follows that \(z^*z\precsim V\).  Proper infiniteness, block
addition, and induction give \(V\precsim v\) for every finite \(r\).
Now \(z^*z\to d\), and the norm-limit rule yields \(d\precsim v\).
This proof includes \(v=0\).
\end{proof}

In particular, Cuntz equivalence preserves proper infiniteness, and
a block sum of properly infinite positives is properly infinite.
These statements follow from block addition and transitivity.
Also, \(u\precsim v\) implies \(\Id(u)\subset\Id(v)\), since
each comparison approximant belongs to \(\Id(v)\).

\begin{lemma}[Transport after a cutoff]\label{lem:transport}
Let \(C\) be a \(C^*\)-algebra, let \(m,n\ge1\) be integers,
and let \(d\in M_n(C)_+\) and \(v\in M_m(C)_+\).
Suppose \(d\precsim v\) in \(C\otimes\Kc\).  For
\(\varepsilon>0\), put
\(h=(d-\varepsilon)_+\).  There is \(w\in M_{m,n}(C^{**})\)
such that
\begin{equation}\label{eq:transport-supports}
 w^*w=e_h,\qquad ww^*\le e_v,
\end{equation}
and
\begin{equation}\label{eq:transport-membership}
 w\Her_{M_n(C)}(h)\subset M_{m,n}(C),\qquad
 w\Her_{M_n(C)}(h)w^*\subset\Her_{M_m(C)}(v).
\end{equation}
Conjugation by \(w\) is an injective \(*\)-homomorphism on
\(\Her(h)\).  For \(l\in\Her(h)_+\),
\begin{equation}\label{eq:transport-cutoffs}
 (wlw^*-\tau)_+=w(l-\tau)_+w^*,\qquad
 w(l-\tau)_+w^*\CuEq(l-\tau)_+\quad(\tau\ge0),
\end{equation}
and
\begin{equation}\label{eq:transport-hereditary-image}
 w\Her(l)w^*=\Her(wlw^*).
\end{equation}
Every representation, with matching matrix amplifications, satisfies
\begin{equation}\label{eq:transport-represented-norm}
 \|\pi^{(m)}(wkw^*)\|=\|\pi^{(n)}(k)\|
 \qquad(k\in\Her(h)).
\end{equation}
The products in these formulas are algebra elements; \(w\) itself
is only asserted to lie in the bidual.
\end{lemma}

\begin{proof}
Compress a stabilization comparison coefficient to the first \(m\)
target and first \(n\) source coordinates.  These scalar multiplier
projections give \(s\in M_{m,n}(C)\) with
\(\|d-s^*vs\|<\varepsilon/2\).  Put \(z=v^{1/2}s h\).
Order in the unitization and functional calculus for \(d\) give
\begin{equation}\label{eq:transport-order}
 \frac{\varepsilon}{2}h^2
 \le z^*z=h s^*vs h
 \le(\|d\|+\varepsilon/2)h^2.
\end{equation}
For the lower bound, use \(h d h\ge\varepsilon h^2\).
The two inequalities and Lemma~\ref{lem:comparison-rules}(i) show
\(\Her(z^*z)=\Her(h)\).  Consequently
\(\Her(|z|)=\Her(h)\) and \(e_{z^*z}=e_h\).
Let \(z=w|z|\) be the rectangular polar decomposition in the
bidual.  Then \(w^*w=e_h\); the factorization through \(v^{1/2}\)
gives \(ww^*=e_{zz^*}\le e_v\).

For \(k\in\Her(h)=\Her(|z|)\), take cutoff functions \(g_j\) for
\(|z|\), zero near zero, with \(g_j(|z|)k\to k\).  Since
\[
 w g_j(|z|)=z\,\frac{g_j(|z|)}{|z|}\in M_{m,n}(C),
\]
norm closure gives \(wk\in M_{m,n}(C)\).  If \(k\ge0\), then
\(wkw^*=(wk^{1/2})(wk^{1/2})^*\) belongs to \(M_m(C)\).
It is supported under \(e_v\), hence belongs to \(\Her(v)\).
Linear decomposition handles arbitrary \(k\).
The identity \(w^*w=e_h\) proves multiplicativity and injectivity.
It also proves \eqref{eq:transport-cutoffs}: functional calculus
vanishing at zero commutes with the homomorphism, and the two
squares of \(w(l-\tau)_+^{1/2}\) are the asserted positives.

For \eqref{eq:transport-hereditary-image}, normal functional
calculus gives \(e_{wlw^*}=we_lw^*\).  Thus the image of
\(\Her(l)\) is supported under \(e_{wlw^*}\), which gives one
inclusion by Lemma~\ref{lem:comparison-rules}(i).  A sandwich
by the image of \(l\) has the form
\[
 (wlw^*)b(wlw^*)
 =w\bigl((wl)^*b(wl)\bigr)w^*.
\]
The inner algebra element is supported on both sides by \(e_l\),
so lies in \(\Her(l)\).  Since conjugation is an injective
\(*\)-homomorphism on \(\Her(l)\), its image is norm closed;
taking closures gives the reverse inclusion.  Finally, the normal
extension of \(\pi\) carries
\eqref{eq:transport-supports} to the same identities with
represented supports.  Conjugating back by the represented \(w^*\)
recovers \(\pi^{(n)}(k)\).  The two contractive conjugations give
both inequalities in \eqref{eq:transport-represented-norm}.
\end{proof}

If \(d_0\in M_\ell(C)_+\) and \(d=d_0^{\oplus r}\in
M_{r\ell}(C)_+\), split \(w\in M_{m,r\ell}(C^{**})\) into
block columns \(w_i\in M_{m,\ell}(C^{**})\).  Then
\[
 w_i^*w_j=\delta_{ij}e_{(d_0-\varepsilon)_+}.
\]
The \(r\) copies have orthogonal ranges and the same initial support.
The single rectangular map also transports every entry of a matrix
over the common hereditary algebra.  Returning to the notation of
Lemma~\ref{lem:transport}, for \(j\ge1\) the same support
identities apply to
\(1_j\otimes w\in M_{jm,jn}(C^{**})\).  Thus, for every
\(K\in M_j(\Her_{M_n(C)}(h))\), entrywise conjugation gives an
algebra matrix and preserves its norm in every representation, with
the matching \(jm\) and \(jn\) amplifications.  This is the joint
norm assertion needed when a finite tuple is transported.

For orthogonal \(c_1,\ldots,c_r\in C_+\), the row
\(R=[c_1^{1/2},\ldots,c_r^{1/2}]\in M_{1,r}(C)\) satisfies
\[
 R^*R=\operatorname{diag}(c_1,\ldots,c_r),\qquad
 RR^*=\sum_i c_i.
\]
Thus their block sum and their sum in \(C\) are Cuntz equivalent.
Moreover, \(f(\sum_i c_i)=\sum_i f(c_i)\) for every continuous
\(f\) with \(f(0)=0\), first for polynomials with zero constant term
and then by continuous functional calculus.  Apply this calculation
in \(M_m(C)\) to the orthogonal transported copies
\(\theta_i(y)=w_i y w_i^*\) of
\(y\in\Her_{M_\ell(C)}((d_0-\varepsilon)_+)_+\).
The cutoff and comparison assertions in
Lemma~\ref{lem:transport} give
\begin{equation}\label{eq:orthogonal-copy-cutoffs}
 \left(\sum_{i=1}^r\theta_i(y)-\tau\right)_+
 \CuEq ((y-\tau)_+)^{\oplus r}\qquad(\tau\ge0).
\end{equation}
In particular this includes the zero cutoff.

\begin{lemma}[Ideals and spectral cutoffs]\label{lem:ideals-cutoffs}
Let \(D\) be a \(C^*\)-algebra.
\begin{enumerate}[label=\textup{(\roman*)}]
\item If \(J=\overline{\bigcup_\lambda J_\lambda}\) for an
increasing directed family of closed ideals and \(z\in J_+\), then
for every \(\varepsilon>0\), \((z-\varepsilon)_+\) lies in some
\(J_\lambda\).  In particular, the ideals
\(\Id((z-\delta)_+)\), directed as \(\delta\downarrow0\), densely
exhaust \(\Id(z)\).
\item For \(u,v\in D_+\), one has \(u\in\Id(v)\) if and only if
every irreducible representation \(\pi\) with \(\pi(v)=0\)
also satisfies \(\pi(u)=0\).
\item If \(f,g\) are continuous nonnegative functions of a
positive element \(a\), vanish at zero, and have the same open
positivity set in \(\operatorname{sp}(a)\), then
\(\Her(f(a))=\Her(g(a))\) and \(f(a)\CuEq g(a)\).
\end{enumerate}
\end{lemma}

\begin{proof}
For (i), choose \(b\in J_\lambda\) with \(\|z-b\|<\varepsilon\).
Then \(\|z+J_\lambda\|<\varepsilon\), and functional calculus in
the quotient gives \((z-\varepsilon)_+\in J_\lambda\).
Since \((z-\delta)_+\to z\), the particular exhaustion follows.
For (ii), apply separation of points by irreducible representations
to the quotient \(D/\Id(v)\).  For (iii), the two support
projections coincide.  Lemma~\ref{lem:comparison-rules}(i) gives
equality of the hereditary algebras and mutual comparison.
\end{proof}

\begin{definition}\label{def:rich}
A positive contraction \(a\) is \emph{rich} if
\((a-\varepsilon)_+\) is properly infinite for every
\(\varepsilon\ge0\).
\end{definition}

Zero is rich, and under \(\propP\) every positive contraction is
rich.  A \emph{ramp} is a continuous nondecreasing function with
values in \([0,1]\), zero on a neighborhood of zero and constant
at \(1\) beyond a specified larger threshold.  A ramp of a rich
contraction is rich.  Indeed, for \(0\le\tau<1\), the open
positivity set of \((f(a)-\tau)_+\) is either empty or the
positivity set of \((a-\rho)_+\) for some \(\rho\ge0\).
Lemma~\ref{lem:ideals-cutoffs}(iii) and richness apply; for
\(\tau\ge1\) the cutoff is zero.  Lemma~\ref{lem:transport}
gives the same conclusion for transported ramps.
